% Vietnamese translation for OI Public Library
% Source-File: IOI中国国家候选队论文/2005/朱晨光/朱晨光.ppt
\documentclass[11pt]{article}
\usepackage{oipl-vn}

\title{Bản trình chiếu: Tư tưởng nhân đôi trong thi tin học}
\author{Zhu Chenguang}
\date{2005}

\begin{document}
\maketitle

\origtitle{Slide companion for doubling methods}
\sourcepath{IOI China candidate team papers / 2005 / Zhu Chenguang / Zhu Chenguang.ppt}

\section{Phạm vi bản dịch}

Bản trình chiếu đi cùng bài viết về tư tưởng nhân đôi. Các slide khôi phục
được chứa công thức phân tích \(n\) theo lũy thừa hai, độ phức tạp
\(O(\log^2 n)\), ví dụ Dijkstra, và bảng thưa.

\section{Lũy thừa và bước nhảy}

Nhân đôi biến một thao tác lặp \(n\) lần thành tổ hợp các bước có độ dài
\(1,2,4,\ldots\). Nếu đã biết kết quả sau \(2^j\) bước, ta có thể ghép hai
kết quả để được \(2^{j+1}\) bước. Đây là nền tảng của binary lifting, lũy
thừa nhanh, và nhiều cấu trúc bảng.

\section{Bảng thưa}

Với bảng thưa, \(B[i,j]\) lưu thông tin trên đoạn bắt đầu tại \(i\) có độ dài
\(2^j\). Sau khi tiền xử lý \(O(N\log N)\), truy vấn RMQ có thể lấy hai đoạn
chồng phủ:
\[
  [i,i+2^k-1],\qquad [j-2^k+1,j],
\]
từ đó trả lời trong \(O(1)\) với phép kết hợp phù hợp.

\section{Kết luận}

Tư tưởng nhân đôi xuất hiện khi quá trình có tính kết hợp hoặc có thể ghép
hai nửa giống nhau. Nó đổi thời gian tuyến tính theo số bước thành thời gian
logarit theo số bước.

\end{document}
